% !TeX program = xelatex
% Compile with local XeLaTeX. Certificate images use absolute filesystem paths.
\documentclass[UTF8,a4paper,11pt,fontset=fandol]{ctexart}
\usepackage[top=18mm,bottom=19mm,left=20mm,right=20mm,headheight=14pt,headsep=5mm]{geometry}
\usepackage{fontspec}
\setmainfont{texgyretermes-regular.otf}[BoldFont=texgyretermes-bold.otf,ItalicFont=texgyretermes-italic.otf]
\setmonofont{lmmono10-regular.otf}[Scale=0.9]
\usepackage{graphicx,booktabs,tabularx,array,float,caption,fancyhdr}
\usepackage[hidelinks,unicode]{hyperref}
\usepackage{hyperref}
\hypersetup{pdftitle={第一周学习周报},pdfauthor={He Song}}
\setlength{\parindent}{2em}
\setlength{\parskip}{3pt}
\setlength{\emergencystretch}{1em}
\setlength{\intextsep}{8pt}
\linespread{1.06}
\raggedbottom
\ctexset{section={format=\large\bfseries,beforeskip=10pt,afterskip=5pt}}
\captionsetup{font=small,labelfont=bf,labelsep=quad,skip=4pt}
\pagestyle{fancy}
\fancyhf{}
\fancyhead[L]{\small Ma's Lab}
\fancyhead[R]{\small 第一周学习周报}
\fancyfoot[C]{\small\thepage}
\renewcommand{\headrulewidth}{0.3pt}

\begin{document}
\begin{center}
{\LARGE\bfseries 第一周学习周报}\par
\vspace{6pt}
{\normalsize He Song\quad\textbar\quad 2026 年 10 月 3 日}
\end{center}
\vspace{2pt}

\noindent\textbf{摘要：}本周期完成了 freeCodeCamp 响应式网页设计课程及 JavaScript 部分，以及 Kaggle Python、Pandas 两门微课程；使用 Git 管理学习项目并推送至 GitHub；并选取了相关数据集以便后续使用。同时订阅了相关学者的学术动态。通过这些学习，我初步了解了网页开发、数据处理与版本管理在后续实践中的作用。目前主要的疑惑是，大模型 Agent 与可视化、HCI 在具体研究中如何结合。

\section{任务要求}
《课题组新成员入门学习计划》将第 1--2 周设为基础技能阶段，主要要求如下。

\begin{table}[H]
\centering
\small
\caption{基础技能阶段的学习要求}
\label{tab:requirements}
\renewcommand{\arraystretch}{1.18}
\begin{tabularx}{\linewidth}{@{}>{\raggedright\arraybackslash\bfseries}p{30mm}X@{}}
\toprule
学习部分 & 应完成的内容 \\
\midrule
Web 开发 & 完成 Responsive Web Design；JavaScript 学至基础算法部分，重点掌握 DOM 操作与 Promise、\texttt{fetch} 等异步编程。 \\
Python 数据处理 & 学习 Python 基础，完成 Kaggle Python 和 Pandas 微课程；第 2 周末选定公开数据集，完成清洗、基本统计及一页数据概况。 \\
Git 与 GitHub & 掌握 git、commit、branch、merge 和 push，完成 Introduction to GitHub 入门练习。 \\
\bottomrule
\end{tabularx}
\end{table}

\section{Web 开发学习}
已完成 freeCodeCamp 响应式网页设计课程并取得认证，JavaScript 的“DOM 操作”和“异步编程”部分也已完成。

通过学习，我更清楚了 HTML、CSS 与 JavaScript 的分工。尤其是描述性 HTML 和 表现性 HTML 的区别是我之前没有涉足的。这些内容为后续制作交互式展示页面打下了基础。

\begin{figure}[H]
\centering
\includegraphics[width=112mm]{/Users/mason/Desktop/Ma's Lab/Getting Start/Cartificates/ResponseWebDesing.png}
\caption{freeCodeCamp 响应式网页设计证书}
\label{fig:web}
\end{figure}

\clearpage
\section{Python 与 pandas 学习}
已完成 Kaggle Python 和 Pandas 两门微课程。我的主要收获是了解了如何用 Python 组织数据处理过程，并用 pandas 进行筛选与统计。

\begin{figure}[H]
\centering
\includegraphics[width=110mm]{/Users/mason/Desktop/Ma's Lab/Getting Start/Cartificates/QianMoA - Python.png}
\caption{Kaggle Python 完成证书}
\label{fig:python}
\end{figure}

\begin{figure}[H]
\centering
\includegraphics[width=110mm]{/Users/mason/Desktop/Ma's Lab/Getting Start/Cartificates/QianMoA - Pandas.png}
\caption{Kaggle Pandas 完成证书}
\label{fig:pandas}
\end{figure}

\section{版本控制与学术关注}
已使用 Git 管理学习项目，并通过 \texttt{git push} 提交到了 GitHub 的 \textit{Private Repo} 中。

已订阅 Huamin Qu、Ashish Vaswani 和 Duen Horng Chau 的学术动态，希望逐步了解相关研究。

\section{数据集的选取}
我选取了数据集 \href{https://www.kaggle.com/datasets/adilshamim8/salaries-for-data-science-jobs}{Data Science, AI \& ML Job Salaries in 2025} 作为我后续任务的数据集。

通过 python 代码，对数据进行了初步的清洗和整理：

python 代码：
\begin{verbatim}
import pandas as pd

url = "/Users/mason/Desktop/Ma's Lab/Getting Start/Dataset/salaries.csv"
raw = pd.read_csv(url)

print("数据规模：", raw.shape)
print(raw.head())
raw.info()

data = raw.copy()
data.columns = data.columns.str.strip()

for col in data.select_dtypes(include=["object", "string"]).columns:
    data[col] = data[col].str.strip().replace("", pd.NA)

data["salary_in_usd"] = pd.to_numeric(
    data["salary_in_usd"], errors="coerce"
)

print("\n各列缺失值：")
print(data.isna().sum())

print("\n完全相同的额外行数：", data.duplicated().sum())

invalid_salary = (
    data["salary_in_usd"].isna()
    | data["salary_in_usd"].le(0)
)
print("无法用于薪资分析的行数：", invalid_salary.sum())

clean = data.loc[~invalid_salary].copy()

print("\n美元薪资描述统计：")
print(clean["salary_in_usd"].describe().round(2))

print("\n记录数最多的 5 个岗位：")
print(clean["job_title"].value_counts().head(5))

print("\n按经验等级统计薪资：")
print(
    clean.groupby("experience_level")["salary_in_usd"]
    .agg(["count", "mean", "median"])
    .reindex(["EN", "MI", "SE", "EX"])
    .round(2)
)

print("\n公司所在地区占比（前 5）：")
print(
    clean["company_location"]
    .value_counts(normalize=True)
    .head(5)
    .mul(100)
    .round(2)
)
\end{verbatim}

输出结果：
\begin{verbatim}

数据规模： (151445, 11)
   work_year experience_level employment_type       job_title  ...  
0       2025               EX              FT    Head of Data  ... 
1       2025               EX              FT    Head of Data  ...
2       2025               SE              FT  Data Scientist  ...
3       2025               SE              FT  Data Scientist  ...                            
4       2025               MI              FT        Engineer  ...

  employee_residence remote_ratio  company_location company_size
0                 US            0                US            M
1                 US            0                US            M
2                 US            0                US            M
3                 US            0                US            M
4                 US          100                US            M

[5 rows x 11 columns]
<class 'pandas.core.frame.DataFrame'>
RangeIndex: 151445 entries, 0 to 151444
Data columns (total 11 columns):
 #   Column              Non-Null Count   Dtype 
---  ------              --------------   ----- 
 0   work_year           151445 non-null  int64 
 1   experience_level    151445 non-null  object
 2   employment_type     151445 non-null  object
 3   job_title           151445 non-null  object
 4   salary              151445 non-null  int64 
 5   salary_currency     151445 non-null  object
 6   salary_in_usd       151445 non-null  int64 
 7   employee_residence  151445 non-null  object
 8   remote_ratio        151445 non-null  int64 
 9   company_location    151445 non-null  object
 10  company_size        151445 non-null  object
dtypes: int64(4), object(7)
memory usage: 12.7+ MB

各列缺失值：
work_year             0
experience_level      0
employment_type       0
job_title             0
salary                0
salary_currency       0
salary_in_usd         0
employee_residence    0
remote_ratio          0
company_location      0
company_size          0
dtype: int64

完全相同的额外行数： 79532
无法用于薪资分析的行数： 0

美元薪资描述统计：
count    151445.00
mean     157527.46
std       74150.77
min       15000.00
25%      105800.00
50%      146100.00
75%      198000.00
max      800000.00
Name: salary_in_usd, dtype: float64

记录数最多的 5 个岗位：
job_title
Data Scientist       18751
Software Engineer    16948
Data Engineer        16352
Data Analyst         13779
Engineer             11004
Name: count, dtype: int64

按经验等级统计薪资：
                  count       mean    median
experience_level                            
EN                13663   98961.74   86600.0
MI                46128  142301.18  130000.0
SE                87491  172657.76  161000.0
EX                 4163  200471.43  189370.0

公司所在地区占比（前 5）：
company_location
US    89.52
CA     4.10
GB     2.76
AU     0.39
NL     0.32
Name: proportion, dtype: float64
\end{verbatim}

\section{疑惑与希望获得的指导}
Agent 主要研究如何让大模型自主完成复杂任务，看起来与可视化、HCI 有些差异，为什么课题组会同时选择这三个方向？“为 Agent 设计人机协作界面”是否是它们的主要连接点，还是还有其他交叉方式？

\end{document}
